\section*{Sets and Indices}

\begin{itemize}
    \item Products \(P=\{X,Y\}\), indexed by \(p\).
\end{itemize}

\section*{Parameters}

All numerical values are read from the CSV files using the field names shown below.

\begin{itemize}
    \item \(a_p\): machine time requirement per batch of product \(p\), in minutes (from \texttt{Machine\_Time\_minutes} in \texttt{table\_1.csv}).
    \item \(b_p\): craftsman time requirement per batch of product \(p\), in minutes (from \texttt{Craftsman\_Time\_minutes} in \texttt{table\_1.csv}).
    \item \(M\): available machine hours (code data: \texttt{machine\_time\_available}).
    \item \(R\): available regular craftsman hours (code data: \texttt{craftsman\_regular\_time\_available}).
    \item \(O\): available overtime craftsman hours (code data: \texttt{craftsman\_overtime\_available}).
    \item \(c^m\): machine-hour cost (code data: \texttt{machine\_time\_cost}).
    \item \(c^r\): regular craftsman-hour cost (code data: \texttt{craftsman\_time\_cost}).
    \item \(c^o\): overtime craftsman-hour cost (code data: \texttt{craftsman\_overtime\_cost}).
    \item \(v_X\): revenue per batch of product \(X\) (code data: \texttt{revenue\_product\_X}).
    \item \(v_Y\): revenue per batch of product \(Y\) (code data: \texttt{revenue\_product\_Y}).
    \item \(L_X\): minimum required batches of product \(X\) (code data: \texttt{min\_batches\_product\_X}).
    \item \(T_Y\): Product \(Y\) QA threshold in batches (code data: \texttt{product\_Y\_QA\_threshold}).
    \item \(F_Y\): one-time QA fee for exceeding the Product \(Y\) threshold (code data: \texttt{product\_Y\_QA\_fee}).
    \item \(\bar M_Y = 10000\): big-\(M\) constant used in the implemented QA-tier constraint.
\end{itemize}

Define the implied resource-use expressions
\[
h^m(x,y)=\frac{a_X x+a_Y y}{60}, \qquad
h^c(x,y)=\frac{b_X x+b_Y y}{60}.
\]

\section*{Decision Variables}

\begin{itemize}
    \item \(x\) (code: \texttt{X}): batches of product \(X\), continuous, \(x \ge 0\).
    \item \(y\) (code: \texttt{Y}): batches of product \(Y\), continuous, \(y \ge 0\).
    \item \(r\) (code: \texttt{craft\_reg}): regular craftsman hours used, continuous, \(0 \le r \le R\).
    \item \(o\) (code: \texttt{craft\_ot}): overtime craftsman hours used, continuous, \(0 \le o \le O\).
    \item \(q\) (code: \texttt{product\_Y\_QA}): binary indicator for the Product \(Y\) QA fee, \(q \in \{0,1\}\).
\end{itemize}

\section*{Objective}

Maximize weekly profit:
\[
\max \;
v_X x + v_Y y
- c^m \, h^m(x,y)
- c^r r
- c^o o
- F_Y q.
\]

Equivalently,
\[
\max \;
v_X x + v_Y y
- c^m \frac{a_X x+a_Y y}{60}
- c^r r
- c^o o
- F_Y q.
\]

\section*{Constraints}

\begin{align}
h^m(x,y) &\le M
&& \text{(machine capacity)} \\
h^c(x,y) &= r+o
&& \text{(craftsman-time balance)} \\
x &\ge L_X
&& \text{(minimum production of product \(X\))} \\
y &\le T_Y + \bar M_Y q
&& \text{(implemented Product \(Y\) QA-tier logic)} \\
0 \le r &\le R
&& \text{(regular craftsman availability)} \\
0 \le o &\le O
&& \text{(overtime craftsman availability)} \\
x &\ge 0,\; y \ge 0,\; q \in \{0,1\}. &&
\end{align}

\section*{Notes on Implementation}

\begin{itemize}
    \item The model is implemented as a mixed-integer linear program and solved with \texttt{GUROBI\_CMD}.
    \item Product quantities \(x\) and \(y\) are continuous, so fractional batches are allowed.
    \item The regular/overtime split is modeled through variables \(r\) and \(o\) with the equality
    \[
    h^c(x,y)=r+o,
    \]
    so all craftsman hours required by production must be assigned to either regular or overtime hours.
    \item The QA revision for Product \(Y\) is implemented exactly as
    \[
    y \le T_Y + 10000\, q,
    \]
    with fee term \(-F_Y q\) in the objective. This formulation allows \(q=1\) both when \(y>T_Y\) and when \(y\le T_Y\); the fee is therefore economically discouraged rather than enforced by a reverse implication. The formulation above matches the code and should not be strengthened when reproducing the workspace model.
    \item There is no explicit upper bound on \(y\) other than resource limits and the big-\(M\) QA-tier constraint.
\end{itemize}
